\section{沿七条问题线确定下一步，而非按新奇图像选题}
第3章的相反结果经常来自不同问题或不同条件，不能一律当作尚未解决的逻辑矛盾。下表把已有证据与可区分的下一步连接起来；“未解决”指本报告已核读证据仍不足，并非断言世界上没有后续研究。
\begin{longtable}{P{25mm}P{55mm}P{66mm}}
\toprule 问题线 & 当前证据留下的具体缺口 & 有区分力的下一步\\\midrule\endfirsthead
\toprule 问题线 & 当前证据留下的具体缺口 & 有区分力的下一步\\\midrule\endhead
快率（\ref{story:fast}） &
色散关联已有反例；层长与供磁仍会随参数调整。 &
保持外层条件可比，分别关闭/改变Hall、惯性和压力响应，同时检查实际X点破冻结预算。\\
全局（\ref{story:global}） &
局地快层能与反弹、旁路和尺寸减速共存。 &
固定宏观形状与初始片宽的两种扫描分别进行；统一累计磁通、上游量和完成时间定义。\\
启动（\ref{story:onset}） &
次级岛可能促进、阻塞或仅伴随重联；递归临界标度未完全统一。 &
测量每代电流片的局地尺度、寿命与扰动，扫描网格和噪声，区分线性增长与净能量释放。\\
闭合（\ref{story:closure}） &
长电子层、热流和三维模态尚不能由一套闭合普遍复现。 &
把F611状态方程与梯度十矩在同一初值比较；检验上游标签、层长、压力散度及能量一致性。\\
第三维（\ref{story:three}） &
有限长X线的受抑尺度主要来自PIC特定序列，Hall/双流体能复现到哪一步未定。 &
共同初边值下扫描质量比、导场、$\beta$和边界；同时验证$B_z$输运、活跃长度、压力项与连接变化。\\
热化（\ref{story:energy}） &
轨道展宽与高温区关联尚未唯一识别不可逆机制。 &
沿粒子轨道与流体体积双重记账；比较中心矩、压力应变、热流、碰撞熵产和诊断平均。\\
观测（\ref{story:diagnostic}） &
嵌套四极和下游升温均有不止一种解释。 &
联合场、流、压力与时间尺度；用多点重建或合成轨道传播坐标及采样误差。\\
\bottomrule
\end{longtable}

本次最直接可由Hall/双流体继续尝试的问题，是\textbf{有限长度抑制与压力闭合是否相互改变}：电子磁通输运能否正确给出活跃/受抑区，而闭合又是否正确给出变薄时间和长层？这是一项由第\ref{story:closure}与第\ref{story:three}节共同引出的研究设计，不是已证实的新规律。若只把PIC压力项换成任意电阻后碰巧得到相近总率，还不足以完成这个检验。\citep{F611,F659,F792}

原书目还缺少若干连接环节，已补登记Ohia2012、Ohia2015和Jain2009，但尚未取得并完成全文核读。它们暂不作为本报告结论证据。原有优先队列中的EMHD撕裂渐近/数值对照、三维电子独有重联、KH跨模型比较和新$\beta$标度也继续保留；不把“列入下一批”写成已经发现或证实其题名所述机制。
